\section{原则三：返回有意义的上下文}

工具的返回值设计同样至关重要。工具实现应该只返回\textbf{高信号}（high signal）的信息给 Agent，优先考虑上下文相关性而非灵活性。

\subsection{避免低级技术标识符}

\begin{table}[H]
\centering
\caption{返回字段设计：低信号 vs 高信号}
\begin{tabular}{p{0.35\textwidth}p{0.35\textwidth}}
\toprule
\textbf{低信号字段（避免）} & \textbf{高信号字段（推荐）} \\
\midrule
\texttt{uuid} & \texttt{name} \\
\texttt{256px\_image\_url} & \texttt{image\_url} \\
\texttt{mime\_type} & \texttt{file\_type} \\
\bottomrule
\end{tabular}
\end{table}

Agent 处理自然语言的名称、术语或标识符的能力\textbf{远优于}处理密码学标识符。Anthropic 发现，仅将任意字母数字 UUID 解析为更具语义意义的可解释语言（甚至是从 0 开始的 ID 方案），就能\textbf{显著提高 Claude 在检索任务中的精确度}，减少幻觉。

\subsection{灵活的响应格式}

在某些场景下，Agent 可能需要同时使用自然语言标识符和技术标识符来触发下游 Tool 调用。例如：

$$\texttt{search\_user(name='jane')} \rightarrow \texttt{send\_message(id=12345)}$$

解决方案是暴露一个简单的 \texttt{response\_format} 枚举参数，让 Agent 自行控制工具返回"简洁"（concise）还是"详细"（detailed）响应：

\begin{lstlisting}[language=Python, caption={ResponseFormat 枚举示例}]
from enum import Enum

class ResponseFormat(Enum):
    DETAILED = "detailed"  # 返回完整信息
    CONCISE = "concise"    # 仅返回必要标识符
\end{lstlisting}

还可以添加更多格式以获得更大灵活性，类似于 GraphQL 中选择性请求特定字段的能力。

\begin{warningbox}{响应结构的选择没有银弹}
工具响应的结构——XML、JSON 还是 Markdown——对评估性能有实际影响。LLM 基于 next-token prediction 训练，倾向于更好地处理与训练数据匹配的格式。没有一种万能的解决方案，最优的响应结构因任务和 Agent 而异，需要通过评估来确定。
\end{warningbox}

\subsection{本章小结}

工具返回值应追求高信号、语义化。避免返回低级技术标识符，优先使用自然语言和可解释的命名。通过 \texttt{response\_format} 参数提供灵活性，并根据评估选择最优的响应结构格式。


